% !TEX program = xelatex
\documentclass[12pt,a4paper]{report}

% ===== 中文與字型 =====
\usepackage{fontspec}
\usepackage{xeCJK}
% 依照課堂範例 week3-1002 的字型設定：標楷體 + SimSun + Times New Roman
% 標楷體沒有粗體，所以用 AutoFakeBold 產生粗體效果
\IfFontExistsTF{標楷體}{%
  \setCJKmainfont[AutoFakeBold=2, AutoFakeSlant=0.2]{標楷體}
  \setCJKsansfont[AutoFakeBold=2, AutoFakeSlant=0.2]{標楷體}
  \setCJKmonofont[AutoFakeBold=2, AutoFakeSlant=0.2]{SimSun}
  \setmainfont{Times New Roman}
}{%
  % 非 Windows 電腦的備用字型
  \setCJKmainfont[AutoFakeBold=2, AutoFakeSlant=0.2]{Noto Serif CJK TC}
  \setCJKsansfont[AutoFakeBold=2, AutoFakeSlant=0.2]{Noto Sans CJK TC}
  \setCJKmonofont[AutoFakeBold=2, AutoFakeSlant=0.2]{Noto Sans CJK TC}
}
% 中文自動斷行（課堂範例的必要設定）
\XeTeXlinebreaklocale "zh"
\XeTeXlinebreakskip = 0pt plus 1pt
\linespread{1.3}

% ===== 版面與套件 =====
\usepackage[margin=2.5cm]{geometry}
\usepackage{xcolor}
\usepackage{graphicx}
\usepackage{booktabs}
\usepackage{tabularx}
\usepackage{enumitem}
\usepackage{tikz}
\usetikzlibrary{arrows.meta,positioning,shapes}
\usepackage{listings}
\usepackage{caption}
\usepackage{titlesec}
% 章標題改為「第 X 章」
\titleformat{\chapter}[display]{\bfseries\huge}{第\,\thechapter\,章}{0.6em}{\Huge}
\titlespacing*{\chapter}{0pt}{0pt}{1.5em}
\usepackage[hidelinks]{hyperref}

% ===== 中文名稱 =====
\renewcommand{\contentsname}{目錄}
\renewcommand{\listfigurename}{圖目錄}
\renewcommand{\listtablename}{表目錄}
\renewcommand{\figurename}{圖}
\renewcommand{\tablename}{表}
\renewcommand{\chaptername}{第}
\renewcommand{\lstlistingname}{程式碼}
\renewcommand{\lstlistlistingname}{程式目錄}
\renewcommand{\abstractname}{摘要}

% ===== listings 設定：左側自動行號 =====
\definecolor{codebg}{RGB}{247,247,247}
\definecolor{kw}{RGB}{0,0,180}
\definecolor{cm}{RGB}{0,120,0}
\definecolor{str}{RGB}{160,30,30}
\lstset{
  basicstyle=\ttfamily\small,
  numbers=left,
  numberstyle=\tiny\color{gray},
  stepnumber=1,
  numbersep=8pt,
  backgroundcolor=\color{codebg},
  frame=single,
  rulecolor=\color{gray},
  breaklines=true,
  breakatwhitespace=false,
  showstringspaces=false,
  tabsize=4,
  captionpos=t,
  keywordstyle=\color{kw}\bfseries,
  commentstyle=\color{cm},
  stringstyle=\color{str},
  extendedchars=false,
  columns=fullflexible,
  keepspaces=true
}

\title{\textbf{LINE Messaging API 註冊申請教學文件}\\[1em]\large 專業排版與設計 期中考}
\author{資管三B \quad 王品翔 \quad 413630715}
\date{\today}

\begin{document}

% ===== 封面 =====
\begin{titlepage}
  \centering
  \vspace*{3cm}
  {\Huge\bfseries LINE Messaging API\\[0.5em] 註冊申請教學文件\par}
  \vspace{1.5cm}
  {\Large 專業排版與設計 期中考\par}
  \vspace{3cm}
  \begin{tabular}{rl}
    班級： & 資管三B \\[0.5em]
    學號： & 413630715 \\[0.5em]
    姓名： & 王品翔 \\
  \end{tabular}
  \vfill
  {\large \today\par}
\end{titlepage}

% ===== 摘要 =====
\begin{abstract}
本文件以繁體中文介紹 LINE Messaging API 的完整申請與設定流程，內容涵蓋 LINE Developers 帳號註冊、Provider 與 Messaging API Channel 的建立、Channel Secret 與 Channel Access Token 的取得與保管、Webhook 的設定與以 ngrok 進行本機測試，最後以 Python 與 Flask 實作一個回聲機器人（Echo Bot）並逐行說明程式碼。讀者依照本文件步驟操作，即可在本機完成第一個可接收並回覆訊息的 LINE Bot。
\end{abstract}

% ===== 四種目錄 =====
\tableofcontents
\listoffigures
\listoftables
\lstlistoflistings

% =========================================================
\chapter{LINE Messaging API 簡介與應用場景}

\section{什麼是 LINE Messaging API}
LINE Messaging API 是 LINE 官方提供給開發者的介面，讓開發者能透過「LINE 官方帳號」與使用者進行雙向溝通。當使用者傳送訊息給官方帳號時，LINE 平台會將事件以 HTTP POST 的方式送到開發者指定的 Webhook 網址；開發者的伺服器處理完後，再呼叫 Messaging API 將回覆訊息送回給使用者。整體運作架構如圖~\ref{fig:arch} 所示。

\begin{figure}[htbp]
  \centering
  \begin{tikzpicture}[
      node distance=2.6cm,
      box/.style={draw, rounded corners, minimum width=2.6cm, minimum height=1.1cm, align=center, fill=blue!8},
      arr/.style={-{Stealth[length=3mm]}, thick}]
    \node[box] (user) {使用者\\LINE App};
    \node[box, right=of user] (line) {LINE\\Platform};
    \node[box, right=of line] (bot) {Bot 伺服器\\(Flask)};
    \draw[arr] ([yshift=3mm]user.east) -- node[above]{\small 傳送訊息} ([yshift=3mm]line.west);
    \draw[arr] ([yshift=3mm]line.east) -- node[above]{\small Webhook} ([yshift=3mm]bot.west);
    \draw[arr] ([yshift=-3mm]bot.west) -- node[below]{\small Reply API} ([yshift=-3mm]line.east);
    \draw[arr] ([yshift=-3mm]line.west) -- node[below]{\small 回覆訊息} ([yshift=-3mm]user.east);
  \end{tikzpicture}
  \caption{LINE Messaging API 運作架構}
  \label{fig:arch}
\end{figure}

\section{常見應用場景}
Messaging API 能應用在許多實際情境中，表~\ref{tab:usecase} 整理了常見的應用類型。

\begin{table}[htbp]
  \centering
  \caption{LINE Bot 常見應用場景}
  \label{tab:usecase}
  \begin{tabularx}{0.9\textwidth}{lX}
    \toprule
    \textbf{應用類型} & \textbf{說明} \\
    \midrule
    客服機器人 & 自動回答常見問題，降低人工客服負擔 \\
    訂單與預約通知 & 以推播訊息通知出貨、預約時間或繳費提醒 \\
    會員服務 & 結合 LIFF 與 Rich Menu 提供點數查詢、會員卡 \\
    資訊查詢 & 查詢天氣、課表、公車到站時間等即時資訊 \\
    行銷推廣 & 透過 Flex Message 推送圖文並茂的優惠活動 \\
    \bottomrule
  \end{tabularx}
\end{table}

\section{主要訊息類型}
Messaging API 提供多種回覆方式，最常用的是 \texttt{Reply Message}（以 reply token 回覆使用者剛傳來的訊息，免費且不計入額度）與 \texttt{Push Message}（主動推播給使用者，會計入每月訊息額度）。此外還有 \texttt{Multicast}、\texttt{Broadcast} 等大量發送方式，以及 Flex Message、Rich Menu 等豐富的版面元件。

% =========================================================
\chapter{LINE Developers 帳號註冊與 Channel 建立}

\section{登入 LINE Developers Console}
\begin{enumerate}[label=步驟 \arabic*：, leftmargin=*]
  \item 開啟瀏覽器前往 \url{https://developers.line.biz/}，點選右上角「Log in to Console」。
  \item 使用個人 LINE 帳號登入（可用 QR Code 或電子郵件登入）。
  \item 第一次登入時需建立開發者帳號，輸入名稱與電子郵件並同意開發者條款。
\end{enumerate}

\section{建立 Provider}
Provider 代表提供服務的個人或組織，一個 Provider 底下可以有多個 Channel。在 Console 首頁點選「Create a new provider」，輸入 Provider 名稱（例如自己的名字或專案名稱）後按「Create」即可。

\section{建立 Messaging API Channel}
目前 LINE 的規定是先建立「LINE 官方帳號」，再於官方帳號中啟用 Messaging API。完整流程如圖~\ref{fig:flow} 所示。

\begin{figure}[htbp]
  \centering
  \begin{tikzpicture}[
      node distance=0.55cm,
      stp/.style={draw, rounded corners, fill=green!8, minimum width=8.5cm, minimum height=0.85cm, align=center},
      arr/.style={-{Stealth[length=2.5mm]}, thick}]
    \node[stp] (a) {1. 登入 LINE Developers Console};
    \node[stp, below=of a] (b) {2. 建立 Provider};
    \node[stp, below=of b] (c) {3. 建立 LINE 官方帳號（LINE Official Account Manager）};
    \node[stp, below=of c] (d) {4. 在官方帳號設定中啟用 Messaging API 並選擇 Provider};
    \node[stp, below=of d] (e) {5. 回到 Console 取得 Channel Secret 與 Access Token};
    \node[stp, below=of e] (f) {6. 設定 Webhook URL 並開啟 Use webhook};
    \foreach \x/\y in {a/b,b/c,c/d,d/e,e/f} \draw[arr] (\x) -- (\y);
  \end{tikzpicture}
  \caption{Messaging API Channel 建立流程}
  \label{fig:flow}
\end{figure}

建立官方帳號時需填寫的欄位整理如表~\ref{tab:fields}。

\begin{table}[htbp]
  \centering
  \caption{建立官方帳號時的必填欄位}
  \label{tab:fields}
  \begin{tabularx}{0.9\textwidth}{lX}
    \toprule
    \textbf{欄位} & \textbf{說明} \\
    \midrule
    帳號名稱 & 顯示給使用者看的 Bot 名稱，之後可修改 \\
    電子郵件 & 用於接收 LINE 官方通知 \\
    公司所在國家或地區 & 選擇「台灣」 \\
    業種（大／小分類） & 依用途選擇，例如「個人」或「教育」 \\
    \bottomrule
  \end{tabularx}
\end{table}

\section{建議的回應設定}
在 LINE Official Account Manager 的「回應設定」中，建議將「聊天」關閉、「Webhook」開啟，並關閉「自動回應訊息」與「加入好友的歡迎訊息」，避免官方預設訊息和自己的程式回覆重複。

% =========================================================
\chapter{Channel Secret 及 Access Token 取得與設定}

\section{Channel Secret}
Channel Secret 位於 Channel 頁面的「Basic settings」分頁。它的用途是驗證 Webhook 請求確實來自 LINE 平台：LINE 會用 Channel Secret 對請求內容做 HMAC-SHA256 運算，並放在 \texttt{X-Line-Signature} 標頭中，伺服器端再以同樣方式比對。

\section{Channel Access Token}
切換到「Messaging API」分頁，捲動到最下方的「Channel access token (long-lived)」，按下「Issue」即可產生長期有效的 Access Token。呼叫 Reply、Push 等 API 時，都必須在 HTTP 標頭中帶上此 Token。

\begin{table}[htbp]
  \centering
  \caption{兩種憑證的比較}
  \label{tab:cred}
  \begin{tabularx}{0.95\textwidth}{lXX}
    \toprule
     & \textbf{Channel Secret} & \textbf{Channel Access Token} \\
    \midrule
    位置 & Basic settings 分頁 & Messaging API 分頁 \\
    用途 & 驗證 Webhook 簽章 & 呼叫 Messaging API 的授權 \\
    使用端 & 接收訊息時 & 送出訊息時 \\
    外洩風險 & 他人可偽造請求 & 他人可冒用 Bot 發訊息 \\
    \bottomrule
  \end{tabularx}
\end{table}

\section{以環境變數保管憑證}
憑證絕對不能直接寫在程式碼裡或上傳到 GitHub。建議在專案目錄建立 \texttt{.env} 檔，並將其加入 \texttt{.gitignore}，內容格式如程式碼~\ref{lst:env}。

\begin{lstlisting}[caption={.env 環境變數設定範例}, label={lst:env}]
LINE_CHANNEL_SECRET=請填入你的ChannelSecret
LINE_CHANNEL_ACCESS_TOKEN=請填入你的ChannelAccessToken
\end{lstlisting}

% =========================================================
\chapter{Webhook 的設定與測試方式}

\section{Webhook 的概念}
Webhook 是 LINE 平台主動「回呼」開發者伺服器的機制。由於 LINE 平台必須能從網際網路連到你的伺服器，且網址必須是 HTTPS，因此在本機開發時，需要透過 ngrok 將本機埠號轉成公開的 HTTPS 網址。

\section{安裝 Python 套件}
首先安裝範例程式需要的套件，指令如程式碼~\ref{lst:pip}。

\begin{lstlisting}[language=bash, caption={安裝所需 Python 套件}, label={lst:pip}]
pip install flask line-bot-sdk python-dotenv
\end{lstlisting}

\section{使用 ngrok 建立公開網址}
\begin{enumerate}[label=步驟 \arabic*：, leftmargin=*]
  \item 至 ngrok 官網註冊帳號並下載 ngrok。
  \item 依官網指示設定 authtoken。
  \item 先啟動 Flask 程式（監聽 8000 埠），再另開一個終端機執行程式碼~\ref{lst:ngrok} 的指令。
  \item 複製畫面上 \texttt{Forwarding} 欄位的 \texttt{https://xxxx.ngrok-free.app} 網址。
\end{enumerate}

\begin{lstlisting}[language=bash, caption={ngrok 啟動指令}, label={lst:ngrok}]
ngrok config add-authtoken <你的authtoken>
python app.py
ngrok http 8000
\end{lstlisting}

\section{在 Console 設定 Webhook URL}
回到 Channel 的「Messaging API」分頁，在「Webhook URL」欄位填入 ngrok 網址並加上路徑 \texttt{/callback}，例如 \texttt{https://xxxx.ngrok-free.app/callback}。接著按「Verify」，若顯示 Success 代表連線成功，最後記得開啟「Use webhook」開關。

\section{常見錯誤排除}
\begin{table}[htbp]
  \centering
  \caption{Webhook 常見錯誤與解決方式}
  \label{tab:error}
  \begin{tabularx}{0.95\textwidth}{lX}
    \toprule
    \textbf{錯誤現象} & \textbf{可能原因與解決方式} \\
    \midrule
    Verify 回傳 404 & 網址少了 \texttt{/callback} 路徑 \\
    Verify 回傳 400 & Channel Secret 錯誤導致簽章驗證失敗 \\
    Verify 逾時 & Flask 未啟動，或 ngrok 埠號與程式不一致 \\
    Bot 沒有回覆 & 未開啟 Use webhook，或 Access Token 錯誤 \\
    重開 ngrok 後失效 & 免費版網址會變動，需重新填寫 Webhook URL \\
    \bottomrule
  \end{tabularx}
\end{table}

% =========================================================
\chapter{範例程式說明：接收及回覆 LINE 訊息}

\section{LINE Bot 範例程式碼}
以下範例使用 Python 與 Flask 框架實作一個簡單的回聲機器人（Echo Bot），會把使用者傳來的文字訊息原封不動地回覆，完整程式碼如程式碼~\ref{lst:app}。

\lstinputlisting[language=Python, caption={Python app.py 範例}, label={lst:app}]{app.py}

\section{程式碼逐段說明}
\begin{itemize}[leftmargin=*]
  \item \textbf{匯入套件（第 1--8 行）：} 匯入 Flask 網頁框架與 LINE Bot SDK v3 的相關類別，並匯入 \texttt{load\_dotenv} 讀取環境變數。
  \item \textbf{讀取憑證（第 10--20 行）：} 從 \texttt{.env} 取得 Channel Secret 與 Access Token，分別建立 \texttt{WebhookHandler}（負責驗證與分派事件）及 \texttt{Configuration}（負責呼叫 API）。
  \item \textbf{callback 函式（第 23--35 行）：} 定義 \texttt{/callback} 路由接收 LINE 的 POST 請求，取出 \texttt{X-Line-Signature} 與請求內容交給 handler 驗證；簽章不符時回傳 HTTP 400，成功則回傳 \texttt{OK}。
  \item \textbf{handle\_message 函式（第 38--47 行）：} 以裝飾器註冊「文字訊息事件」的處理函式，利用事件中的 \texttt{reply\_token} 建立 \texttt{ReplyMessageRequest}，將使用者的文字原樣回覆。
  \item \textbf{啟動伺服器（第 50--51 行）：} 在 0.0.0.0 的 8000 埠啟動 Flask，搭配 ngrok 即可對外提供服務。
\end{itemize}

\section{測試結果}
在 Console 的「Messaging API」分頁掃描 QR Code 將 Bot 加為好友，傳送任何文字，例如「你好」，Bot 就會立刻回覆「你好」，代表整個流程已成功運作。

\section{延伸應用}
完成 Echo Bot 後，可以進一步嘗試以下功能：使用 \texttt{push\_message} 主動推播通知；建立 Rich Menu 提供常用功能按鈕；或以 Flex Message 設計卡片式的圖文訊息，讓 Bot 的介面更豐富。

% =========================================================
\chapter{結論}
本文件從 LINE Messaging API 的基本概念出發，依序說明開發者帳號註冊、Channel 建立、憑證取得與保管、Webhook 設定與 ngrok 測試，最後以 Python 實作並解析一個完整可運作的 Echo Bot。掌握這些步驟後，開發者即可在此基礎上延伸出客服、通知、查詢等各式實用的 LINE 應用。

% =========================================================
\chapter{本文件 LaTeX 原始碼}
本章列出產生本文件的完整 \LaTeX{} 原始碼。

\lstinputlisting[language={[LaTeX]TeX}, caption={本文件 LaTeX 原始碼 main.tex}, label={lst:self}]{main.tex}

\end{document}
